%-------------------------------------------------------------------%
% Pengaturan tambahan untuk Templat Laporan TA Bentuk 2
% (Proyek Kelompok bersama Mitra / Capstone), Prodi IF ITERA.
% Dipanggil dari thesis.tex setelah config/config.sty.
%-------------------------------------------------------------------%
\usepackage{colortbl}
\usepackage{amssymb}
\usepackage{pdfpages}

%----- Judul bab: satu baris, tengah, tebal, huruf kapital ---------%
% Contoh: "BAB I MASALAH MITRA DAN TUJUAN PROYEK"
\titleformat{\chapter}[block]%
  {\large\centering\bfseries}%
  {\MakeUppercase{\chaptertitlename}\ \thechapter}{0.4em}%
  {\MakeUppercase}
\titlespacing*{\chapter}{0pt}{0pt}{12pt}

%----- Nomor halaman: tengah bawah pada semua halaman --------------%
\fancypagestyle{capstone}{%
  \fancyhf{}%
  \renewcommand{\headrulewidth}{0pt}%
  \fancyfoot[C]{\thepage}%
}
\fancypagestyle{plain}{%
  \fancyhf{}%
  \renewcommand{\headrulewidth}{0pt}%
  \fancyfoot[C]{\thepage}%
}

%----- Saklar petunjuk ---------------------------------------------%
% \petunjuktrue  = teks "Petunjuk:" dan halaman Petunjuk Pengisian
%                  serta Peta Kewajiban DITAMPILKAN (saat menulis).
% \petunjukfalse = semua petunjuk DISEMBUNYIKAN (naskah final).
\newif\ifpetunjuk
\petunjuktrue

%----- Penanda status kewajiban ------------------------------------%
\newcommand{\wajib}[1]{\textcolor{red!75!black}{\textbf{\upshape #1}}}
\newcommand{\opsional}[1]{\colorbox{yellow!55}{\textbf{\upshape #1}}}
\newcommand{\tidakperlu}[1]{\textcolor{gray}{\upshape #1}}
% [Seminar proposal: ... · Sidang akhir: ...]
\newcommand{\tahap}[2]{[Seminar proposal: #1 $\cdot$ Sidang akhir: #2]}

%----- Kotak petunjuk (miring, abu-abu, spasi 1) -------------------%
\newcommand{\petunjuk}[1]{%
  \ifpetunjuk
    \par\begingroup\setlength{\parindent}{0pt}\singlespacing\footnotesize\itshape
    \textcolor{gray!60!black}{\textbf{Petunjuk:} #1}\par\endgroup\vspace{4pt}%
  \fi}

%----- Teks isian (abu-abu dalam kurung siku) ----------------------%
\newcommand{\isian}[1]{\textcolor{gray}{[#1]}}
\newcommand{\tulisuraian}{\par\textcolor{gray}{[Tulis uraian dalam bentuk paragraf.]}\par}

%----- Judul halaman tanpa nomor yang masuk daftar isi -------------%
% Contoh: \judulhalaman{KATA PENGANTAR}
\newcommand{\judulhalaman}[1]{%
  \chapter*{#1}%
  \phantomsection\addcontentsline{toc}{chapter}{#1}%
  \thispagestyle{capstone}}

%----- Subbab tanpa nomor otomatis (P.3.1, dst.) --------------------%
\newcommand{\subbabP}[1]{%
  \section*{#1}%
  \phantomsection\addcontentsline{toc}{section}{#1}}

%----- Halaman pembatas BAGIAN I / BAGIAN II -----------------------%
\newcommand{\halamanbagian}[2]{%
  \chapter*{#1}%
  \phantomsection\addcontentsline{toc}{chapter}{#1}%
  \thispagestyle{capstone}%
  {\centering\itshape #2\par}}

%----- Sel tabel berwarna untuk Peta Kewajiban ---------------------%
\newcommand{\cW}[1]{\cellcolor{red!12}\textcolor{red!75!black}{\textbf{#1}}}
\newcommand{\cO}[1]{\cellcolor{yellow!35}\textbf{#1}}
\newcommand{\cT}[1]{\cellcolor{gray!15}\textcolor{gray!80!black}{#1}}
\newcommand{\cekbox}{$\square$}

%----- Variabel informasi laporan ----------------------------------%
\makeatletter
\newcommand{\proyekTim}[1]{\def\@proyektim{#1}}
\newcommand{\printproyektim}{\@proyektim}
\newcommand{\anggotaLain}[1]{\def\@anggotalain{#1}}
\newcommand{\printanggotalain}{\@anggotalain}
\newcommand{\namaMitra}[1]{\def\@namamitra{#1}}
\newcommand{\printnamamitra}{\@namamitra}
\newcommand{\tempatLahir}[1]{\def\@tempatlahir{#1}}
\newcommand{\printtempatlahir}{\@tempatlahir}
\newcommand{\tahunLaporan}[1]{\def\@tahunlaporan{#1}}
\newcommand{\printtahunlaporan}{\@tahunlaporan}
\newcommand{\tanggalLaporan}[1]{\def\@tanggallaporan{#1}}
\newcommand{\printtanggallaporan}{\@tanggallaporan}
\newcommand{\koordinatorProdi}[2]{\def\@namakoor{#1}\def\@nipkoor{#2}}
\newcommand{\printnamakoor}{\@namakoor}
\newcommand{\printnipkoor}{\@nipkoor}
\makeatother
